\documentclass[letterpaper]{article} % DO NOT CHANGE THIS
\usepackage[submission]{aaai2027}  % DO NOT CHANGE THIS
\usepackage[hyphens]{url}  % DO NOT CHANGE THIS
\usepackage{graphicx} % DO NOT CHANGE THIS
\urlstyle{rm} % DO NOT CHANGE THIS
\def\UrlFont{\rm}  % DO NOT CHANGE THIS
\usepackage{natbib}  % DO NOT CHANGE THIS AND DO NOT ADD ANY OPTIONS TO IT
\usepackage{caption} % DO NOT CHANGE THIS AND DO NOT ADD ANY OPTIONS TO IT
\frenchspacing  % DO NOT CHANGE THIS

\usepackage{booktabs}
\usepackage{multirow}
\usepackage{amsmath}
\usepackage{xcolor}

\pdfinfo{/TemplateVersion (2027.1)}

\setcounter{secnumdepth}{1}

\title{Can Small Language Models Serve Ordinary Users? A Lower-Bound Benchmark under Non-Expert Prompts and Edge Constraints}

\author{Anonymous Submission}
\affiliations{Anonymous Institution}

\begin{document}
\maketitle

\begin{abstract}
Small language models are increasingly promoted as practical alternatives to cloud-scale LLMs for users with limited hardware access, privacy concerns, or edge deployment needs. Yet most evaluations assume idealized conditions: expert-written prompts, clean benchmark inputs, and full-precision inference on powerful infrastructure. In practice, ordinary users often provide ambiguous or poorly structured prompts, while edge deployment commonly relies on compressed or quantized models. This paper studies whether small language models remain reliable when both user input and deployment conditions are imperfect. We introduce \textbf{PEEL} (\textbf{P}rompt--\textbf{E}dge \textbf{E}valuation \textbf{L}adder), an evaluation protocol for measuring model reliability as prompt quality and deployment conditions degrade. PEEL varies prompt quality, from expert-formatted to colloquially broken prompts, and deployment quality, from full-precision inference to aggressive low-bit quantization. We evaluate 18 small and efficient language models across reading comprehension, open-domain question answering, truthfulness discrimination, instruction following, and multi-turn dialogue. Our results show that quantization alone is not the dominant source of failure once output-scoring artefacts are handled. Instead, prompt degradation causes broader and more severe capability loss, especially in span-based reading comprehension. We further find that architecture affects robustness under imperfect prompting, with hybrid efficient models showing stronger instruction-following performance than comparable small baselines. PEEL highlights the gap between clean benchmark assumptions and ordinary-user edge conditions, and provides a practical protocol for lower-bound evaluation of small language models.
\end{abstract}


\section{Introduction}
\label{sec:intro}

Large language models (LLMs) have rapidly expanded the capabilities of natural language processing, enabling general-purpose assistance across reasoning, question answering, summarization, coding, and dialogue. However, these gains have largely been driven by increasing model scale, training data, and computational resources. As a result, many state-of-the-art LLMs remain difficult to deploy outside cloud-based infrastructure, where inference can be costly, latency-sensitive applications may suffer, and user data must often be transmitted to remote servers. These limitations are particularly important for privacy-sensitive or resource-constrained scenarios, where local or edge deployment is preferred.

Small language models (SLMs) have therefore emerged as a practical path toward more accessible and efficient AI. With fewer parameters and lower memory requirements, SLMs can reduce inference cost and support deployment on low-power or edge devices. Prior work highlights that such models are especially attractive for resource-limited settings, where on-device processing can improve privacy, response time, personalization, and cloud independence. In this sense, SLMs are not merely smaller versions of LLMs; they are a potential mechanism for broadening access to language technologies under realistic hardware and connectivity constraints.

However, the accessibility promise of SLMs remains incomplete. Existing evaluations often measure models under idealized conditions: prompts are carefully written by experts, benchmark inputs are clean and well-structured, and deployment is performed using full-precision or well-supported inference settings. These conditions reflect an upper-bound view of model capability. They do not necessarily reflect how ordinary users interact with language models in practice. Many users do not have prompt engineering expertise; they may provide vague, incomplete, colloquial, or poorly formatted instructions. At the same time, they may rely on compressed or quantized models due to limited hardware, cost, privacy, or connectivity constraints.

This creates a neglected reliability problem. If SLMs are expected to make AI more accessible, then their evaluation should account for the conditions faced by the users they are meant to serve. A model that performs well with expert prompts on full-precision cloud infrastructure may fail when deployed as a quantized model and queried with non-expert prompts. Conversely, some models may preserve useful capabilities under these imperfect conditions. We argue that model evaluation should therefore include not only upper-bound benchmark performance, but also lower-bound reliability: how well models behave when both user input and deployment conditions degrade.

% Large language models are increasingly deployed beyond cloud datacenters: on
% laptops, smartphones, and embedded devices with strict memory and compute
% budgets. Simultaneously, AI adoption is broadening to users with no machine
% learning background, who interact with models through informal, under-specified,
% or syntactically broken natural-language prompts. These two trends jointly
% define a \emph{real-world lower-bound scenario} that current benchmarks were
% not designed to measure.

% Most public benchmarks measure an \emph{upper bound}. They evaluate carefully
% formatted prompts, assume the model is accessed through a high-quality serving
% stack, and often treat decoding, chat templates, and output parsing as
% implementation details. That convention is useful for comparing frontier
% capabilities, but it hides a practical question faced by edge deployments:
% \emph{what happens when a small model is both converted for local inference and
% queried by a non-expert user?} A 4-bit checkpoint may preserve benchmark
% accuracy under a canonical prompt, yet fail when the same task is written as a
% short casual request. Conversely, an informal prompt may be harmless for a
% cloud-served chat model but brittle when the runtime no longer applies the
% model's chat template.

% Prior benchmarks each cover part of this space. PalmBench~\citep{palmbench2024}
% and HELM~\citep{helm2022} offer broad capability coverage, but largely assume
% structured inputs and production-quality inference. RobustAlpacaEval~\citep{ralpaca2024}
% and PromptBench~\citep{promptbench2023} study prompt perturbations, but do not
% cross them with deployment conversion. Quantization work such as GPTQ and
% GGUF/\texttt{llama.cpp}~\citep{gptq2022,gguf} makes local inference practical,
% but usually reports accuracy under expert prompts. As a result, the field has
% strong evidence for \emph{how quantization behaves when prompting is fixed} and
% for \emph{how prompting behaves when serving is fixed}, but limited evidence for
% their interaction.

% This interaction matters because failures can be nonlinear. Edge conversion
% changes not only weight precision, but also the runtime interface: chat-template
% insertion, stop-token handling, tokenizer compatibility, and output post-processing
% can all differ from the reference HuggingFace path. Low-quality prompts stress
% the same interface because they omit explicit answer formats and rely on the
% model to infer the task. The user-visible lower bound is therefore a joint
% property of model architecture, quantization level, runtime, prompt quality, and
% scoring protocol. EdgeBench is designed to make that joint property measurable.

% \paragraph{Contributions.}
% \begin{enumerate}
%   \item A \textbf{2D lower-bound evaluation framework} crossing five edge
%     deployment settings with four prompt quality tiers.
%   \item A \textbf{four-tier prompt taxonomy} (P0--P3) with hand-crafted
%     templates for reading comprehension, open-domain QA, and truthfulness
%     discrimination.
%   \item A \textbf{reproducible evaluation suite} spanning 18 models,
%     1,184 reference-scored runs, and 150 MT-Bench judge-scored runs across
%     20 experimental conditions.
%   \item Empirical evidence of \textbf{task-type-dependent degradation}: SQuAD
%     EM remains near 20--22\% under P0 but collapses to 0\% under P3,
%     contrasting with TruthfulQA's 4.8\,pp setting-level change, and
%     \textbf{architecture-dependent prompt
%     sensitivity}: Falcon-H1 (SSM-Hybrid) models are significantly more robust to informal
%     prompting than same-size Transformers at instruction-following tasks.
% \end{enumerate}

\section{Related Work}
\label{sec:related}

\paragraph{LLM Benchmarks.}
HELM~\citep{helm2022} and BIG-Bench~\citep{bigbench2022} establish multitask
evaluation at scale by collecting diverse tasks, standardizing metrics, and
making model comparisons reproducible. MMLU~\citep{mmlu2021} and
SuperGLUE~\citep{superglue2019} remain widely used capability probes because
they provide stable reference metrics across knowledge, reasoning, and language
understanding tasks. These benchmarks are intentionally controlled: they remove
many variables that would otherwise confound model comparison. EdgeBench keeps
that reproducibility goal, but changes the target condition. Instead of asking
which model performs best under ideal benchmark prompts, we ask which conclusions
survive when the same task is exposed to edge runtime conversion and informal
prompting.

PalmBench~\citep{palmbench2024} is closest in spirit because it considers
on-device deployment. However, it primarily evaluates capability under a fixed
prompting protocol. HELM-style evaluation similarly offers broad task coverage
but does not treat prompt literacy as an experimental axis. \citet{multiprompt2024}
argue that single-prompt evaluation is unstable and propose multi-prompt
strategies, but their goal is more reliable capability estimation rather than a
lower-bound stress test. EdgeBench differs by deliberately pairing a deployment
axis with a prompt-quality axis, so that a model's reported score is not only a
property of weights and task, but also of the way the model is likely to be used.

\paragraph{Quantization Effects on Accuracy.}
GPTQ~\citep{gptq2022} and GGUF/\texttt{llama.cpp}~\citep{gguf} achieve
practical 4-bit deployment with modest accuracy loss on standard benchmarks.
SqueezeLLM~\citep{squeezellm2023} and QuIP\#~\citep{quip2024} push further.
\citet{edgellm2025} evaluate energy and accuracy of quantized LLMs on ARM
edge platforms; \citet{quantsys2025} provide a systematic characterization
across quantization strategies and bit-widths. Both studies hold prompt quality
fixed at expert level, leaving the interaction with informal user prompts
unstudied.

The quantization literature has also shown that the implementation path matters.
The same nominal bit width can behave differently depending on calibration,
kernel implementation, tokenizer handling, and whether the serving stack inserts
the same system and chat wrappers used during instruction tuning. These details
are often treated as engineering choices outside the benchmark, yet our results
show that they can determine whether a score is interpreted as model weakness or
as a parser/runtime artefact. In particular, the SQuAD and HaluEval corrections
reported later are not new model capabilities; they are examples of how local
inference output formatting can invalidate naive exact-match and yes/no parsers.
This motivates reporting deployment settings as first-class experimental
conditions rather than as undocumented preprocessing.

\paragraph{Prompt Robustness.}
RobustAlpacaEval~\citep{ralpaca2024} measures how model rankings shift under
paraphrase. PromptBench~\citep{promptbench2023} applies adversarial character-
and word-level perturbations to expose fragility. Neither study couples prompt
robustness to edge deployment constraints.

Prompt robustness work typically studies perturbations that are either semantic
paraphrases or adversarial edits. EdgeBench instead uses a user-skill framing:
P0 is a conventional benchmark prompt, P1 is a casual but clear request, P2
removes task framing, and P3 is colloquial and weakly delimited. This framing is
not intended to replace adversarial robustness benchmarks. It captures a
different failure source: ordinary users often omit output constraints, mix
context and question delimiters, or assume the model will infer the task without
being told. Those omissions interact strongly with reference metrics. For
example, SQuAD v2 requires the model to either produce an exact span or the
literal unanswerable marker; removing that instruction changes the task from
span extraction to open-ended conversation. The P0--P3 axis therefore measures a
practical robustness notion that standard paraphrase tests do not isolate.

\paragraph{Novel Architectures.}
State-space models (Mamba~\citep{mamba2023}, Falcon-H1~\citep{falconh12025})
and sub-1-bit BitNet~\citep{bitnet2024} represent emerging families optimized
for efficient inference. \citet{mamba_empirical2024} show Mamba matches
Transformers on many tasks but underperforms on knowledge-intensive benchmarks;
\citet{hymba2024} demonstrate that hybrid SSM+attention architectures outperform
same-size Transformers on standard NLP tasks. EdgeBench is the first study to
measure their prompt-quality robustness and edge quantization tolerance jointly.

These architecture families are especially relevant for edge benchmarking
because the design goal is not only lower memory, but also favorable behavior
under constrained decoding. Falcon-H1 combines Mamba-style selective state
modules with attention layers, while Falcon-E represents a native low-bit model
family rather than a post-hoc quantized dense Transformer. Standard model
leaderboards can understate why these distinctions matter: two models with
similar parameter counts and similar reference accuracy may respond very
differently to a broken prompt or to a runtime that omits chat wrappers. Our
MT-Bench results make this point concrete: Falcon-H1-1.5B substantially
outperforms Qwen2.5-1.5B at E0/P0 and remains stable under P3, while Qwen loses
more than half a point on the same judge-scored instruction-following benchmark.
We therefore treat architecture family as a central moderator, not merely as a
model-table descriptor.

\paragraph{Safety and Alignment for Small Edge Models.}
Safety evaluation is commonly performed on cloud-served chat models whose
alignment layers, system prompts, and refusal behavior are part of a managed
product. Edge deployments weaken this assumption. A local model may be used
without a system prompt, converted to a different runtime, or quantized in a way
that changes refusal phrasing. EdgeBench does not claim to provide a full
red-teaming benchmark; our safety component is a diagnostic refusal analysis
over safe and unsafe prompts. Its purpose is to test whether deployment
conversion preserves a basic safety/helpfulness trade-off. The answer in our
current data is mixed: unsafe refusal is far from reliable, while safe-prompt
over-refusal is non-trivial. This supports a broader lesson shared with the rest
of the benchmark: deployment evaluation should include behavioral checks, not
only task accuracy.

% ─────────────────────────────────────────────────────────────────────────────
\section{Benchmark Design}
\label{sec:design}

\subsection{The 2D Degradation Space}

EdgeBench evaluates every combination of a five-level \emph{edge deployment
axis} and a four-level \emph{prompt quality axis}, forming the 4$\times$5 grid
shown in Figure~\ref{fig:space}.

The grid is intentionally small enough to run exhaustively but structured enough
to separate three effects that are often conflated. First, E0$\to$E1 isolates
GPU-side precision reduction while retaining the HuggingFace inference stack.
Second, E1$\to$E2 crosses from the HuggingFace/chat-template path into a
CPU-oriented GGUF runtime. Third, E2$\to$E4 varies local quantization strength
inside the same runtime family. On the prompt axis, P0$\to$P3 progressively
removes task scaffolding. This lets us distinguish a true numerical
quantization effect from a serving-interface effect and from a user-prompting
effect.

\begin{figure}[t]
  \centering
  \includegraphics[width=\columnwidth]{figures/fig1_benchmark_space.pdf}
  \caption{The EdgeBench 2D evaluation space. Each cell represents one
    experimental condition; colour encodes schematic relative performance drop
    from the top-left (expert prompt, full-precision) baseline.}
  \label{fig:space}
\end{figure}

\paragraph{Edge Deployment Settings (E0--E4).}
\begin{itemize}
  \item \textbf{E0-REF}: BF16, HuggingFace \texttt{transformers}, GPU.
  \item \textbf{E1-INT8}: 8-bit static quantization via \texttt{bitsandbytes}, GPU.
  \item \textbf{E2-Q8}: Q8\_0 GGUF via \texttt{llama.cpp}, CPU.
  \item \textbf{E3-Q5}: Q5\_K\_M GGUF via \texttt{llama.cpp}, CPU.
  \item \textbf{E4-Q4}: Q4\_K\_M GGUF via \texttt{llama.cpp}, CPU.
\end{itemize}
E2--E4 use CPU-only inference to simulate consumer edge hardware with no
discrete GPU.

We treat these settings as \emph{deployment conditions}, not only as bit widths.
This distinction is important because the edge setting changes several
variables at once: kernel implementation, decoding defaults, chat-template
insertion, stop-token handling, and available memory. Reporting them as named
conditions makes the benchmark more faithful to how practitioners deploy local
models, while the E0/E1 and E2/E3/E4 contrasts still provide useful controlled
subcomparisons.

\paragraph{Prompt Quality Tiers (P0--P3).}
We define four tiers aligned with realistic user archetypes drawn from
AI-literacy and human-computer interaction literature~\citep{sclar2023,multiprompt2024}.
The tiers span the spectrum from expert benchmark-style prompting to the informal
interaction style documented in studies of first-time LLM users:
\begin{itemize}
  \item \textbf{P0 -- Expert}: Structured, instruction-following prompt
    identical to standard benchmark format. Explicit task description, labelled
    fields, output format constraint.
  \item \textbf{P1 -- Casual}: Simplified but grammatically correct. Removes
    technical terminology and output format instructions.
  \item \textbf{P2 -- Naive}: Minimal content only; no task framing; relies on
    the model to infer the task type from context alone.
  \item \textbf{P3 -- Broken}: Colloquial, abbreviated, syntactically informal.
    Typical of a first-time LLM user with no prompting background.
\end{itemize}

\paragraph{Prompt-tier internal validity.}
Each tier is defined by a \emph{functional specification} (what information
is withheld and how grammar is degraded) rather than a single template, allowing
future work to instantiate the taxonomy with alternative phrasings.
We validate tier ordering by checking MT-Bench P0--P3 deltas on matched
model--setting pairs. The aggregate effect is small, but the deltas vary
systematically by family: Falcon-H1-1.5B improves slightly from P0 to P3
(8.08$\to$8.20 averaged over settings), while Qwen2.5-1.5B drops from
4.17 to 3.60. This supports treating P1--P3 as a diagnostic stress test,
not as a calibrated user-study scale.
A controlled user study eliciting naturalistic informal prompts and mapping
them to P1--P3 would further validate tier realism; this is flagged as
future work.

The tiers are deliberately task-preserving: every prompt still contains the
question and any required context/options. We do not inject adversarial content,
contradict the original task, or remove the information needed to answer.
Therefore, when performance collapses under P3, the collapse indicates failure
to infer the intended interaction contract rather than absence of task evidence.
This design choice is especially important for SQuAD v2, where the prompt must
communicate both span extraction and the unanswerable convention.

Table~\ref{tab:prompts} illustrates all four tiers for the SQuAD v2 task.

\begin{table}[t]
  \small
  \centering
  \caption{Prompt templates for SQuAD v2 across quality tiers.
    \texttt{\{ctx\}} and \texttt{\{q\}} are filled at run time.}
  \label{tab:prompts}
  \begin{tabular}{@{}lp{5.5cm}@{}}
    \toprule
    Tier & Template \\
    \midrule
    P0 & \textit{Answer the question using only the context. If the answer is
         not in the context, say `unanswerable'. Context: \{ctx\}
         Question: \{q\} Answer:} \\[4pt]
    P1 & \textit{Read the following text and answer the question. If the text
         doesn't have the answer, say so. Text: \{ctx\} Q: \{q\} A:} \\[4pt]
    P2 & \textit{\{q\} $\backslash$n$\backslash$n Here is some info: \{ctx\}} \\[4pt]
    P3 & \textit{pls help, \{q\}?? info: \{ctx\}} \\
    \bottomrule
  \end{tabular}
\end{table}

\subsection{Datasets}

We select tasks that expose different output contracts. SQuAD v2 requires
extractive spans and a special unanswerable marker; NQ-Open requires factual
recall without context; TruthfulQA-MC1 requires selecting among explicit
options; HaluEval requires binary judgment; and MT-Bench evaluates open-ended
multi-turn instruction following. The mix lets us test whether degradation is a
general capability loss or a task-format-specific failure.

\paragraph{Reading Comprehension -- SQuAD v2~\citep{squad2}.}
Extractive QA with unanswerable questions, scored by Exact Match (EM) and F1.
Unanswerable questions require outputting the literal string
\textit{``unanswerable''}, making format adherence part of the evaluation signal.

\paragraph{Open-Domain QA -- Natural Questions Open~\citep{nqopen}.}
Short-answer factual QA without a supplied context passage, scored by EM
against a set of acceptable answer strings. Tests parametric knowledge retention
under prompt simplification.

\paragraph{Truthfulness Discrimination -- TruthfulQA-MC1~\citep{truthfulqa}.}
Single-correct multiple-choice selection from shuffled options, scored by
letter accuracy. We use a generative MC1 protocol (model selects a letter)
rather than log-likelihood scoring, making evaluation fully model-agnostic and
compatible with all runners.

HaluEval~\citep{halueval2023} is included at P0 for hallucination detection
baselines; the judge-style instruction is intrinsic to its prompt and is not
tier-degradable. MT-Bench is run at P0 and P3 because its open-ended prompts
are naturally instruction-following tasks rather than reference-scored QA.

\subsection{Evaluation Protocol}

Reference-scored tasks use fully offline metrics. EM and F1 follow the SQuAD
normalisation convention (lowercase, punctuation stripping). We use 200
randomly sampled examples per (dataset, split) pair, with a fixed seed for
reproducibility; MT-Bench uses all 80 questions and GPT-4o-mini judging.
Total inference cost is dominated by E0/E1 GPU runs; E2--E4 CPU runs require
no GPU allocation.

For aggregate heatmaps, each cell is the mean over available model--setting runs
after filtering reference-scored runs with fewer than 100 examples. MT-Bench is
reported separately because each run contains 80 questions by definition and
would otherwise be incorrectly excluded by the reference-task filter. We report
simple means rather than fitting a hierarchical model because the main paper's
goal is diagnostic: reveal which cells fail, which model families remain stable,
and which scoring artefacts must be corrected before drawing claims. Appendix
tables provide the coverage counts needed to interpret missing or partially
covered cells.

We also audit every result before including it in the claims matrix. The audit
checks that the score is derived from raw model outputs, that the parser matches
the task's output contract, and that missing runs are due to documented runtime
or tokenizer failures rather than selective reporting. This audit led to two
important corrections: SQuAD exact-match scores were recomputed after stripping
EOS/stop tokens, and HaluEval binary scores were recomputed using the first
non-empty output line.

% ─────────────────────────────────────────────────────────────────────────────
\section{Models and Experimental Setup}
\label{sec:models}

\subsection{Model Selection}

We evaluate 18 models spanning nine brand families (Falcon-H1, Falcon-E,
Gemma, InternLM, Llama, Phi, Qwen, SmolLM, TinyLlama) across the
0.09B--3.8B parameter range (Table~\ref{tab:models}).

The model set is organized around three comparison roles. First, Falcon-H1 is
the primary family under study because it represents a recent SSM-Hybrid edge
model line with sizes from 90M to 1.5B parameters. Including 0.5B, 0.6B, 1.5B,
and a deeper 1.5B variant lets us separate scale effects from architecture
effects within the same family. Second, Falcon-E provides native low-bit BitNet
models, which test whether an architecture trained for low precision behaves
differently from post-hoc quantized Transformers. Third, Qwen, Llama, SmolLM,
TinyLlama, Gemma, Phi, and InternLM serve as external baselines that reviewers
will recognize from small-model and edge-LLM literature. The comparison is not
intended to crown a universal winner; it asks whether the same deployment and
prompt perturbations produce consistent failure patterns across families.

We prioritize the 0.5B--3B band because it is the realistic local-deployment
range for consumer laptops and small edge devices. Larger models may achieve
higher absolute scores, but they answer a different systems question: how to
serve capable LLMs with more memory and accelerator support. EdgeBench instead
focuses on the regime where memory pressure makes quantization necessary and
where users are likely to run models through lightweight local runtimes.

\begin{table}[t]
  \small
  \centering
  \caption{Models evaluated in EdgeBench.
    $^\dagger$BitNet 1.58-bit native; excluded from E1--E4 comparisons (no GGUF).
    $^\ddagger$SSM-Hybrid (Mamba-2 + attention layers).}
  \label{tab:models}
  \begin{tabular}{@{}llc@{}}
    \toprule
    Model & Family & Params \\
    \midrule
    Falcon-E-1B$^\dagger$         & Falcon-E    & 1.0B \\
    Falcon-E-3B$^\dagger$         & Falcon-E    & 3.0B \\
    Falcon-H1-90M$^\ddagger$      & Falcon-H1   & 0.09B \\
    Falcon-H1-0.5B$^\ddagger$     & Falcon-H1   & 0.5B \\
    Falcon-H1-0.6B$^\ddagger$     & Falcon-H1   & 0.6B \\
    Falcon-H1-1.5B$^\ddagger$     & Falcon-H1   & 1.5B \\
    Falcon-H1-1.5B-Deep$^\ddagger$& Falcon-H1   & 1.5B \\
    Gemma-3-1B                    & Gemma       & 1.0B \\
    InternLM2.5-1.8B              & InternLM    & 1.8B \\
    Llama-3.2-1B                  & Llama       & 1.0B \\
    Llama-3.2-3B                  & Llama       & 3.0B \\
    Phi-3.5-Mini                  & Phi         & 3.8B \\
    Phi-4-Mini                    & Phi         & 3.8B \\
    Qwen2.5-0.5B                  & Qwen        & 0.5B \\
    Qwen2.5-1.5B                  & Qwen        & 1.5B \\
    SmolLM2-360M                  & SmolLM      & 0.36B \\
    SmolLM2-1.7B                  & SmolLM      & 1.7B \\
    TinyLlama-1.1B                & TinyLlama   & 1.1B \\
    \bottomrule
  \end{tabular}
\end{table}

\subsection{Deployment Configurations}

GPU settings (E0, E1) run on an NVIDIA A100 SXM4 (40\,GB) via Deakin
University's HPC cluster.
CPU settings (E2--E4) use \texttt{llama.cpp} with 4 threads and up to
8\,GB RAM allocation, simulating a mid-range laptop or developer workstation
without a discrete GPU.
BitNet (Falcon-E) models run via \texttt{bitnet.cpp} at native 1.58-bit
precision and are excluded from the E1 INT8 and E2--E4 GGUF tracks.
All experiments ran on the HPC cluster using SLURM job arrays;
full scripts, prompts, and model outputs are released with the paper.

The HPC setting is used for reproducibility and queue management, not to claim
that E2--E4 are high-end server deployments. We pin CPU thread counts and memory
budgets so that the GGUF runs approximate a constrained local environment even
when executed on shared infrastructure. To check that this abstraction is not
misleading, we additionally ran a consumer-workstation replication for selected
Falcon-H1 and Qwen models; the corrected SQuAD scores and throughput trends are
consistent with the HPC GGUF results (Section~\ref{sec:discussion}). This
replication is not used to inflate the main aggregate counts, but it supports
the interpretation of E2--E4 as edge-style conditions.

% ─────────────────────────────────────────────────────────────────────────────
\section{Results}
\label{sec:results}

\subsection{TruthfulQA-MC1: Quantization-Resilient}
\label{sec:tqa}

\begin{figure}[t]
  \centering
  \includegraphics[width=\columnwidth]{figures/fig2_tqa_heatmap.pdf}
  \caption{TruthfulQA-MC1 mean accuracy (\%) across the full 4$\times$5 evaluation grid
    (left) and delta relative to E0-REF (right).
    Chance baseline: 25.0\%. Both axes show modest effects---no systematic collapse.}
  \label{fig:heatmap}
\end{figure}

Across all 1,184 qualifying runs (18 models, num\_examples $\geq$ 100),
TruthfulQA-MC1 mean accuracy at the expert baseline (P0, E0-REF) is 37.9\%;
it rises slightly at E1-INT8 (39.4\%) before settling at 38.4--40.1\% at
E2--E4. All cells remain well above the 25.0\% chance baseline.
Crucially, the P0$\to$P3 drop within any setting is $<$12\,pp, and
quantization from E0 to E4 changes P0 accuracy by only $-$0.7\,pp---a
\emph{near-zero interaction} between the two degradation axes for
this task type.

One notable pattern (Table~\ref{tab:tqa_grid}): at E2--E4 (CPU/llama.cpp),
the P1 ``casual'' prompt tier scores \emph{lower} than at E0/E1
(29.1--31.6\% vs.\ 40.4--41.0\%), while P0 is unaffected.
We attribute this to llama.cpp's lack of a system-level chat template:
casual prompts that omit explicit format instructions rely on the model's
instruction-following mode to infer the task, which is not activated by
llama.cpp's plain-text interface.

This pattern is useful because it separates two hypotheses. If quantization
itself were the dominant cause, P0 should degrade monotonically from E0 to E4.
Instead, P0 remains stable and the largest degradation appears when prompts
remove explicit selection instructions. The result suggests that multiple-choice
truthfulness is numerically robust for small models, but still sensitive to
whether the runtime preserves instruction-following context. In practice, this
means a developer can often quantize MC-style tasks aggressively, but should keep
the prompt format explicit after conversion.

\begin{table}[t]
  \small\centering
  \caption{TruthfulQA-MC1 mean accuracy (\%) in the full 4$\times$5 grid.
    Random-chance baseline: 25.0\%.}
  \label{tab:tqa_grid}
  \begin{tabular}{@{}lcccc@{}}
    \toprule
    Setting & P0 & P1 & P2 & P3 \\
    \midrule
    E0-REF  & 37.9 & 40.4 & 36.7 & 36.6 \\
    E1-INT8 & 39.4 & 41.0 & 35.7 & 37.3 \\
    E2-Q8   & 40.1 & 29.1 & 30.1 & 28.8 \\
    E3-Q5   & 38.4 & 31.6 & 33.0 & 30.1 \\
    E4-Q4   & 38.6 & 30.2 & 32.9 & 30.8 \\
    \bottomrule
  \end{tabular}
\end{table}

\subsection{SQuAD v2: Prompt-Driven Collapse}
\label{sec:collapse}

\begin{figure}[t]
  \centering
  \includegraphics[width=\columnwidth]{figures/fig1_squad_heatmap.pdf}
  \caption{SQuAD v2 Exact Match (\%) across the 4$\times$5 grid (left) and
    delta relative to E0-REF (right, red\,=\,degradation).
    The P3 column collapses to 0\% universally; the P0 row is stable at 20--22\%
    across all five edge settings.}
  \label{fig:collapse}
\end{figure}

After correcting for an EOS/stop-token scoring artefact (Section~\ref{sec:discussion}),
SQuAD v2 EM at P0 is stable at 19.8--22.0\% across all five edge settings,
with no significant capability cliff at the E1$\to$E2 (\texttt{llama.cpp})
boundary (Table~\ref{tab:squad_grid}).
The \emph{dominant failure mode is prompt degradation}, not quantization:
SQuAD EM drops to 0.0\% at P3 for \emph{all} models and settings.
P1 and P2 show gradual degradation (5.8--12.4\% and $\approx$0\%),
but the cliff occurs between P2 and P3.

SQuAD is the clearest example of why lower-bound evaluation cannot be summarized
by an average over prompts. Under P0, the edge settings are surprisingly close:
Q4 and Q5 do not destroy span extraction once stop-token artefacts are removed.
Under P3, the same models receive the same passage and question but no longer
receive a clear extraction contract; they answer conversationally, explain, or
fail to emit the exact span. Exact match then collapses even though the evidence
needed to answer is still present in the context. This is not a generic
``small models are bad'' result. It is a task-contract result: extractive QA
requires explicit formatting more than MC selection does.

\begin{table}[t]
  \small\centering
  \caption{SQuAD v2 Exact Match (\%) in the 4$\times$5 grid.
    P0 is stable across all quantization levels; P3 is universally zero.}
  \label{tab:squad_grid}
  \begin{tabular}{@{}lcccc@{}}
    \toprule
    Setting & P0 & P1 & P2 & P3 \\
    \midrule
    E0-REF  & 19.8 & 7.9 & 0.1 & 0.0 \\
    E1-INT8 & 20.5 & 5.8 & 0.0 & 0.0 \\
    E2-Q8   & 22.0 & 12.4 & 0.2 & 0.0 \\
    E3-Q5   & 21.8 & 10.9 & 0.1 & 0.0 \\
    E4-Q4   & 20.1 & 10.4 & 0.1 & 0.0 \\
    \bottomrule
  \end{tabular}
\end{table}

\subsection{NQ-Open: Quantization Collapse at E2--E4}
\label{sec:nq}

Natural Questions Open reveals a qualitatively different failure mode from SQuAD.
Even at E0-REF, mean EM is only 1.3\% (rising to 1.7\% at E1-INT8)---
reflecting the weak parametric factual recall of sub-2B models without a
supporting context passage. At E2--E4 (llama.cpp CPU), EM drops to
\textbf{0.0\%} across \emph{all} prompt tiers (Table~\ref{tab:nq_grid}).

\begin{table}[t]
  \small\centering
  \caption{NQ-Open Exact Match (\%) in the 4$\times$5 grid.
    E2--E4 are uniformly zero: the E1$\to$E2 boundary marks a complete factual
    recall collapse for sub-2B models on llama.cpp without a chat template.}
  \label{tab:nq_grid}
  \begin{tabular}{@{}lcccc@{}}
    \toprule
    Setting & P0 & P1 & P2 & P3 \\
    \midrule
    E0-REF  & 1.3 & 0.0 & 0.0 & 0.0 \\
    E1-INT8 & 1.7 & 0.0 & 0.0 & 0.0 \\
    E2-Q8   & 0.0 & 0.0 & 0.0 & 0.0 \\
    E3-Q5   & 0.0 & 0.0 & 0.0 & 0.0 \\
    E4-Q4   & 0.0 & 0.0 & 0.0 & 0.0 \\
    \bottomrule
  \end{tabular}
\end{table}

The E1$\to$E2 collapse is caused by two compounding factors:
(1)~\texttt{llama.cpp} does not apply the model's chat template, so models
run in raw-completion mode---NQ-Open requires answer recall without a supporting
passage, which demands active instruction-following; and
(2)~sub-2B models have inherently limited parametric factual recall,
which disappears entirely once instruction-following mode is lost.
This finding underscores a key distinction between
knowledge-intensive open-domain QA and reading comprehension:
quantization matters for the former even under expert prompting.

The NQ result should be read conservatively. Exact match is harsh for open-domain
QA, and the absolute E0/E1 scores show that most evaluated edge-scale models do
not store enough factual knowledge to answer many NQ questions without retrieval.
Nevertheless, the all-zero E2--E4 block is still informative because it identifies
a deployment boundary where even the small amount of recall observed at E0/E1
disappears. For edge applications, this argues against using small local models
as standalone factual QA systems; retrieval augmentation or task-specific
closed-book fine-tuning would be required before NQ-style workloads are
realistic.

\subsection{HaluEval: Binary Classification and Parser Artefacts}
\label{sec:halueval}

HaluEval provides a binary hallucination-detection counterpoint to SQuAD and
NQ-Open. Under the corrected first-line parser, HaluEval-General decreases
from 57.1\% at E0-REF to 43.5\% at E4-Q4, while HaluEval-QA remains in the
41--47\% range across settings (Table~\ref{tab:halueval_grid}). These results
show degradation, but not the complete P3-style collapse observed for SQuAD.

\begin{table}[t]
  \small\centering
  \caption{HaluEval accuracy (\%) under P0 after first-line output parsing.
    The corrected parser removes prompt echoes that otherwise contain both
    ``yes'' and ``no'' and produce ambiguous labels.}
  \label{tab:halueval_grid}
  \begin{tabular}{@{}lccccc@{}}
    \toprule
    Split & E0-REF & E1-INT8 & E2-Q8 & E3-Q5 & E4-Q4 \\
    \midrule
    General & 57.1 & 51.9 & 50.8 & 46.5 & 43.5 \\
    QA      & 41.1 & 39.6 & 47.3 & 47.2 & 43.9 \\
    \bottomrule
  \end{tabular}
\end{table}

The uncorrected parser marked 50--62\% of E2--E4 responses as ambiguous
because verbose outputs often repeated the task instruction after the first
answer. Parsing only the first output line recovers 47--51\% accuracy at
E2--E4, making HaluEval the second example of a llama.cpp-specific evaluation
artefact rather than a pure capability loss.

HaluEval also clarifies the boundary of EdgeBench's conclusions. The corrected
scores are not high in absolute terms; they remain close to a weak binary
classifier. The important result is methodological: without output-aware parsing,
local-runtime evaluations can convert a mediocre but usable answer into an
ambiguous label. This same issue can affect any benchmark whose prompt contains
the label words inside the instruction. We therefore treat parser specification
as part of the benchmark protocol rather than as disposable implementation code.

\subsection{MT-Bench: Instruction Following Under Degradation}
\label{sec:mtbench}

We evaluate MT-Bench (80 multi-turn questions, 8 categories) with GPT-4o-mini
as judge at all five edge settings and two prompt tiers (P0, P3), totalling
150 scored runs. Coverage is not perfectly rectangular: Falcon-E is excluded
from MT-Bench because it uses a separate BitNet runtime, and InternLM/Phi have
partial setting coverage due to tokenizer and chat-template failures.

\begin{figure}[t]
  \centering
  \includegraphics[width=\columnwidth]{figures/fig3_mtbench_family_bar.pdf}
  \caption{MT-Bench scores at E0-REF/P0 (solid) and P3 (transparent) per model,
    colored by family. Falcon-H1 and Gemma cluster at 5--8; most transformer
    baselines fall below 5. Falcon-E is excluded because it uses BitNet runtime
    rather than the HF/GGUF MT-Bench path.}
  \label{fig:mtbench_bar}
\end{figure}

Table~\ref{tab:mtbench_settings} shows mean scores by setting.
Scores range from 4.64 (E4-Q4, P0) to 5.12 (E1-INT8, P0), a
modest 0.49-point spread, confirming that instruction-following \emph{quality}
is largely resilient to quantization---the E3/E4 drop is real but small.
The P0$\to$P3 drop is at most 0.26 points at any single setting.

This MT-Bench pattern contrasts with SQuAD. Open-ended instruction following is
less dependent on exact output form, so informal prompts usually reduce score
gradually rather than catastrophically. The judge can give partial credit for a
reasonable answer even if the model does not follow a strict extraction format.
Thus, MT-Bench is a better measure of conversational usefulness, while SQuAD is
a sharper probe of format adherence. Reporting both prevents a misleading single
story: the same model can be robust conversationally and brittle under exact
reference scoring.

\begin{table}[t]
  \small\centering
  \caption{MT-Bench mean score (1--10) by setting and prompt tier.
    Counts vary by setting because some optional models have partial coverage.}
  \label{tab:mtbench_settings}
  \begin{tabular}{@{}lcc@{}}
    \toprule
    Setting & P0 & P3 \\
    \midrule
    E0-REF  & 5.00 & 4.87 \\
    E1-INT8 & 5.12 & 4.86 \\
    E2-Q8   & 5.02 & 4.92 \\
    E3-Q5   & 4.79 & 4.58 \\
    E4-Q4   & 4.64 & 4.54 \\
    \bottomrule
  \end{tabular}
\end{table}

\subsection{Architecture Family Comparison on MT-Bench}
\label{sec:arch}

Figure~\ref{fig:mtbench_bar} reveals a bimodal distribution.
At E0-REF/P0, Falcon-H1 and Gemma cluster at high scores, with Gemma at 7.43
and the strongest Falcon-H1 variants above 8,
while all other families score below 5.5.
At the matched 1.5B scale, Falcon-H1-1.5B scores \textbf{8.20}
vs.\ Qwen-2.5-1.5B at \textbf{4.06}---a gap of 4.14 points.

The prompt-sensitivity gap is equally striking.
Across the full matched grid, Falcon-H1 models lose only $\mathbf{0.11}$ points
from P0 to P3 as a family average, and the flagship Falcon-H1-1.5B improves
slightly (8.08$\to$8.20 averaged over settings; +0.075 at E0-REF).
In contrast, Qwen2.5-1.5B drops from 4.17 to 3.60 ($-$0.57), and the Llama/
TinyLlama group drops by 0.22 on average.
The finding suggests that Falcon-H1's hybrid selective-state mechanism
provides \emph{structural robustness} to informal prompt framing.

We phrase this as evidence of architecture-dependent robustness rather than a
proof of mechanism. The experiment does not ablate Falcon-H1's internal modules,
so it cannot show that the SSM component alone causes the gain. What it does
show is that, under the same tasks, judge, prompt tiers, and deployment settings,
the Falcon-H1 family occupies a different region of the MT-Bench distribution
from most dense Transformer baselines. That is the relevant benchmark-level
claim: architecture family changes the lower-bound behavior observed by users.

\begin{table}[t]
  \small\centering
  \caption{MT-Bench mean score by family and P0$\to$P3 delta.
    Falcon-E is excluded; InternLM and Phi have partial setting coverage.}
  \label{tab:mtbench_arch}
  \begin{tabular}{@{}lccc@{}}
    \toprule
    Family & P0 & P3 & $\Delta$ (P0$\to$P3) \\
    \midrule
    Gemma       & 7.40 & 7.29 & $-$0.11 \\
    Falcon-H1   & 6.45 & 6.34 & $-$0.11 \\
    InternLM    & 6.78 & 6.69 & $-$0.09 \\
    Llama/TinyLlama & 3.73 & 3.51 & $-$0.22 \\
    Phi         & 3.74 & 3.73 & $-$0.01 \\
    Qwen        & 3.61 & 3.34 & $-$0.27 \\
    SmolLM      & 3.09 & 2.90 & $-$0.19 \\
    \bottomrule
  \end{tabular}
\end{table}

\subsection{Latency vs.\ Accuracy Trade-off}

Table~\ref{tab:latency} reports mean decode throughput by setting.
E0-REF (BF16 GPU) achieves the highest throughput at 116.7\,chars/s.
Counterintuitively, E1-INT8 (GPU int8) is slower at 35.3\,chars/s due to
\texttt{bitsandbytes} dequantization overhead. E2--E4 CPU settings achieve
23.8--31.7\,chars/s without any GPU---practically sufficient for non-real-time
applications. E4-Q4 reduces memory footprint by ${\sim}4\times$ relative to
E0-REF BF16 while retaining ${\sim}$87\% of TruthfulQA accuracy at P0.

\begin{table}[t]
  \small\centering
  \caption{Mean decode throughput (chars/s) by edge setting.}
  \label{tab:latency}
  \begin{tabular}{@{}lccc@{}}
    \toprule
    Setting & Mean & Min & Max \\
    \midrule
    E0-REF (BF16, GPU) & 116.7 & 3.1 & 327.8 \\
    E1-INT8 (int8, GPU) & 35.3 & 0.9 & 89.8 \\
    E2-Q8 (Q8, CPU) & 23.8 & 1.4 & 180.0 \\
    E3-Q5 (Q5\_K\_M, CPU) & 27.0 & 1.1 & 193.7 \\
    E4-Q4 (Q4\_K\_M, CPU) & 31.7 & 1.2 & 232.2 \\
    \bottomrule
  \end{tabular}
\end{table}

\subsection{Safety Alignment Under Edge Deployment}
\label{sec:safety}

We additionally evaluate safety behaviour with HarmBench-style unsafe prompts
and safe control prompts. The current keyword-based classifier assigns each
response to \textsc{Refuse}, \textsc{Comply}, or \textsc{Ambiguous}. Across
75 model--setting pairs, models refuse only 47.3\% of unsafe prompts while
also over-refusing 19.1\% of safe prompts. This indicates that small edge
models are neither reliably safe nor reliably helpful under our safety suite.

\begin{figure}[t]
  \centering
  \includegraphics[width=\columnwidth]{figures/fig5_safety_heatmap.pdf}
  \caption{Safety refusal behaviour across model--setting pairs. Edge
    deployment does not induce a consistent safety trend: some models become
    more compliant under quantization, while others over-refuse benign prompts.}
  \label{fig:safety}
\end{figure}

The most concerning result is not a uniform quantization effect, but instability:
the maximum model--setting refusal rate reaches 84.4\%, whereas several
model--setting pairs remain broadly compliant. Safety behaviour should therefore
be evaluated as a deployment property rather than inferred from full-precision
model identity.

This result also cautions against treating refusal rate alone as alignment
quality. A high refusal rate may indicate safer behavior on unsafe prompts, but
it may also indicate over-refusal on benign prompts. Conversely, a low refusal
rate may reflect helpfulness on safe prompts or unsafe compliance. For this
reason, we report unsafe-prompt refusal and safe-prompt over-refusal separately
and keep the classifier in the appendix as a diagnostic tool rather than a final
safety benchmark.

% ─────────────────────────────────────────────────────────────────────────────
\section{Discussion}
\label{sec:discussion}

\paragraph{What EdgeBench's lower bound means.}
EdgeBench should not be read as a replacement for upper-bound benchmarks. A
model that performs poorly under P3 may still be useful when embedded in a
product with a carefully engineered prompt, retrieval layer, and output parser.
Likewise, a model that is robust under EdgeBench may still fail on tasks outside
our English QA and instruction-following suites. The lower-bound interpretation
is more specific: if a model, runtime, and prompt tier combination fails here,
then a practitioner should not assume that standard benchmark scores will
transfer to a lightly engineered local deployment.

\paragraph{Task type and prompt degradation as primary moderators.}
Our corrected results show that \emph{prompt quality}, not
quantization level, is the dominant driver of SQuAD performance collapse.
SQuAD v2 EM is stable at $\sim$21\% across E0--E4 under expert prompts
(P0), but collapses to $\sim$0\% at P3 for every architecture and quantization
setting. TruthfulQA-MC1 shows the opposite pattern: robust to both prompt
degradation and quantization (4.8\% absolute drop E0$\to$E4).
This task-format interaction motivates separate lower-bound reporting rather
than a single scalar headline metric.

The implication is that benchmark design should preserve task format as a first
class factor. Multiple-choice selection, span extraction, binary judgment, and
open-ended conversation do not degrade in the same way. A single average score
would hide the most actionable lesson in our data: quantized models can remain
usable when the answer space is constrained, while exact-reference tasks become
fragile when the prompt no longer specifies the contract.

\paragraph{Two scoring artefacts and their lesson (W3).}
Our audit identified two independent scoring bugs, both caused by the same
root cause: \texttt{llama.cpp} models at E2--E4 generate verbose, multi-line
outputs containing EOS/stop-token strings and repeated prompt instructions
after the actual answer.

\textit{SQuAD v2 (span EM).}
An initial analysis reported 0\% EM at E2$+$.
The bug: \texttt{[end of text]}, \texttt{<|endoftext|>}, and variant tokens
survive the standard \texttt{normalize\_answer} step, adding extra words that
prevent exact-match against the gold span.
Fix: truncate to the first non-empty output line and strip EOS tokens before
scoring. This recovers SQuAD EM to 21--31\% at E2--E4, consistent with E0-REF.

\textit{HaluEval (yes/no classification).}
A separate audit found that 50--62\% of HaluEval responses at E2--E4 were
marked \emph{ambiguous} (both ``yes'' and ``no'' detected) by the full-text
parser. The cause: quantized models re-generate prompt instructions after
their yes/no answer; these instructions contain both words, so the binary
detector returns ``'' (neither).
Fix: parse only the first output line.
This recovers HaluEval accuracy from $\sim$24\% to $\sim$47--51\% at E2--E4,
revealing that the models' hallucination-detection capability is substantially
better than the uncorrected scores suggested.

\textbf{Both are the same pitfall at a different output level:}
practitioners using \texttt{llama.cpp} at Q8/Q4 who parse free-text outputs
without first-line truncation will systematically underreport accuracy on both
span-extraction \emph{and} binary classification tasks.

This lesson generalizes beyond our tasks. Local runtimes often expose raw text
that a hosted chat API would hide behind message-role formatting and stop-token
logic. Benchmarks that reuse hosted-model parsers on local outputs may therefore
measure interface mismatch rather than capability. We recommend that future
edge-LLM benchmarks publish parsers together with raw outputs and include a
parser audit before reporting model rankings.

\textbf{Consumer-hardware replication.}
To validate the corrected scores, we ran a supplementary local evaluation on an
Intel i7-13700KF workstation (8 threads, CPU-only \texttt{llama.cpp} b8937)
covering Falcon-H1 (our primary SSM-Hybrid family) and Qwen2.5 (Transformer
reference). Using the fixed scorer, \textbf{Falcon-H1-0.5B} achieved 14--16\%
EM (Q8\_0/Q4) and \textbf{Falcon-H1-1.5B} achieved 22\% EM (Q4)---consistent
with the corrected HPC scores (31\%/30.5\% for the same models at E2--E4, with
differences attributable to output verbosity).
Qwen2.5-0.5B similarly reached 20--26\% EM locally vs.\ 24--30\% corrected
HPC. Throughput on the consumer CPU was 29--94 chars/s for Falcon-H1
(vs.\ 24--32 chars/s on HPC), with Q4 faster than Q8 due to memory bandwidth.

\paragraph{Practical recommendations.}
For practitioners deploying models under resource constraints: (1)~prefer
multiple-choice or classification outputs over extractive spans when running
\texttt{llama.cpp} Q4/Q5; (2)~Falcon-H1 SSM-Hybrid models offer better
instruction-following quality and prompt robustness than same-size dense
Transformers at no additional memory cost; and (3)~safety refusal behaviour
must be tested after deployment conversion, because quantization does not
produce a monotonic safety trend.

For benchmark builders, the recommendations are different. First, report
prompt-tier matched results rather than mixing prompt variants into one mean.
Second, keep MT-Bench-style judge scoring separate from reference-scored tasks
because the number of examples, output contract, and scoring noise differ.
Third, document missing cells explicitly. In an edge benchmark, missingness is
itself informative: tokenizer failures, chat-template incompatibilities, and
runtime conversion errors are part of the deployment burden.

\paragraph{Limitations.}
All datasets are English-only. P1--P3 prompt templates are handcrafted;
a controlled user study validating tier realism is future work (see below).
The E1$\to$E2 boundary confound has been partially isolated via a consumer-hardware
replication (Section~Discussion), but a full GGUF-on-GPU ablation within the HPC
cluster remains future work.
Finally, the current safety analysis uses keyword-based refusal detection. It is
useful for detecting instability across model--setting pairs, but it should not
be treated as a substitute for human red teaming or policy-specific safety
evaluation.

% ─────────────────────────────────────────────────────────────────────────────
\section{Conclusion}
\label{sec:conclusion}

We introduced EdgeBench, a benchmark for simultaneously probing LLM performance
along two real-world degradation axes: edge quantization and prompt quality.
Evaluating 18 models across 1,184 reference-scored runs and 150 MT-Bench
judge-scored runs reveals four principal findings:
(1)~task type is the strongest determinant of robustness---MC tasks degrade
gracefully while extraction tasks collapse under informal prompts;
(2)~Falcon-H1-1.5B substantially outperforms same-size Transformers on
instruction-following (8.20 vs.\ 4.06 MT-Bench at 1.5B scale) and remains
stable from P0 to P3;
(3)~quantization-induced output artefacts can masquerade as capability collapse
unless scoring strips stop tokens and prompt echoes; and
(4)~safety refusal is unstable across model--setting pairs, with only 47.3\%
unsafe-prompt refusal and 19.1\% safe-prompt over-refusal in our current
classifier.
We release all prompts, evaluation scripts, and model outputs to support
reproducibility and future lower-bound benchmarking research.

% ─────────────────────────────────────────────────────────────────────────────
\bibliography{references}

\clearpage
\appendix

\section{Appendix: Additional Experimental Details}
\label{app:details}

\subsection{Prompt Tier Operationalization}

The prompt tiers in EdgeBench are intended as a controlled diagnostic axis, not
as a calibrated psychometric model of user skill. P0 preserves the conventional
benchmark prompt with explicit task framing and output constraints. P1 removes
some technical phrasing while keeping grammatical task intent. P2 gives the
model the minimum content needed to infer the task. P3 introduces informal
phrasing, abbreviations, and weak formatting. The current paper uses one
template per task and tier; future work should sample multiple templates per
tier or collect natural prompts from users.

\begin{table*}[t]
  \small\centering
  \caption{Prompt-tier definitions used in the reference-scored tasks.}
  \label{tab:app_prompt_tiers}
  \begin{tabular}{@{}lp{2.6cm}p{9.2cm}@{}}
    \toprule
    Tier & Intended user type & Operational rule \\
    \midrule
    P0 & Benchmark/expert & Explicit instruction, labelled fields, and output-format constraint. \\
    P1 & Casual but clear & Natural wording with task intent preserved; fewer labels and weaker format constraint. \\
    P2 & Naive/minimal & Question and context/options only; no explicit task description. \\
    P3 & Informal/noisy & Casual spelling, abbreviations, and weak delimiters; no formal answer instruction. \\
    \bottomrule
  \end{tabular}
\end{table*}

\subsection{Run Coverage}

Table~\ref{tab:app_coverage} summarizes the coverage used by the main analysis.
Reference-scored tasks use 200 examples per run after filtering; MT-Bench uses
80 questions and is therefore reported separately from the 100-example filter
used for QA/classification aggregates.

\begin{table*}[t]
  \small\centering
  \caption{Run coverage used in the current draft.}
  \label{tab:app_coverage}
  \begin{tabular}{@{}lp{3.0cm}p{8.0cm}@{}}
    \toprule
    Component & Runs & Notes \\
    \midrule
    Reference-scored tasks & 1,184 & 200 examples per run after filtering \\
    MT-Bench & 150 & 80 questions; P0/P3 only \\
    IFEval & 45 & Reported in released artifacts, not central in main text \\
    Safety & 75 model--setting pairs & Keyword/refusal classifier over safe and unsafe prompts \\
    \bottomrule
  \end{tabular}
\end{table*}

\subsection{MT-Bench Coverage and Exclusions}

MT-Bench is evaluated only at P0 and P3 because it is an open-ended
instruction-following benchmark rather than a reference-scored QA task. This
means MT-Bench is not included in the P0/P1/P2/P3 matched heatmaps. Instead,
the paper reports matched P3--P0 deltas by model, family, and setting.
Falcon-E is excluded because it runs through a native BitNet inference path
rather than the HuggingFace/GGUF path used for MT-Bench generation. InternLM
and Phi have partial setting coverage because of tokenizer and chat-template
failures.

\begin{table}[t]
  \small\centering
  \caption{MT-Bench setting-level means.}
  \label{tab:app_mtbench_setting}
  \begin{tabular}{@{}lcc@{}}
    \toprule
    Setting & P0 & P3 \\
    \midrule
    E0-REF  & 5.00 & 4.87 \\
    E1-INT8 & 5.12 & 4.86 \\
    E2-Q8   & 5.02 & 4.92 \\
    E3-Q5   & 4.79 & 4.58 \\
    E4-Q4   & 4.64 & 4.54 \\
    \bottomrule
  \end{tabular}
\end{table}

\subsection{Scoring Artefact Corrections}

Two scoring artefacts were discovered during result auditing. First,
\texttt{llama.cpp} generations for SQuAD often included stop tokens, prompt
echoes, or trailing explanatory text after the answer. Exact-match scoring
therefore underreported span accuracy. The corrected scorer strips known EOS
tokens and truncates to the first non-empty output line before applying standard
SQuAD normalization. Second, HaluEval generations sometimes repeated the
instruction text after the answer; because the instruction contains both
``yes'' and ``no'', a full-text parser marked many outputs as ambiguous. The
corrected parser again uses the first non-empty line.

These corrections are conservative: they do not search the entire response for
a favorable answer, and they do not change the gold labels. They only remove
post-answer text that would not be shown as the answer in a deployed QA or
classification interface.

\subsection{Safety Classifier Caveats}

The safety analysis is based on a keyword/refusal classifier with three labels:
\textsc{Refuse}, \textsc{Comply}, and \textsc{Ambiguous}. The current draft
reports refusal rates rather than attack success rates because ambiguous
responses are common in small edge models. For unsafe prompts, a refusal rate
of 47.3\% means that more than half of unsafe requests are either complied with
or ambiguous. For safe prompts, the 19.1\% refusal rate measures over-refusal.
These numbers should be treated as a first safety diagnostic rather than a full
red-teaming evaluation.

\subsection{Failure Case Taxonomy}

During result auditing, we grouped failures into four recurring types. This
taxonomy is useful for interpreting EdgeBench because not all low scores imply
the same kind of weakness.

\begin{table*}[t]
  \small\centering
  \caption{Failure categories observed during EdgeBench auditing.}
  \label{tab:app_failure_taxonomy}
  \begin{tabular}{@{}p{3.0cm}p{5.6cm}p{5.4cm}@{}}
    \toprule
    Failure type & Observable symptom & Main affected tasks \\
    \midrule
    Task-contract failure & Model answers conversationally, explains, or emits extra text instead of the required span/letter/label. & SQuAD v2, TruthfulQA-MC1 \\
    Runtime-interface failure & Output contains stop tokens, repeated instructions, missing chat-template behavior, or tokenizer-specific artefacts. & SQuAD v2, HaluEval, NQ-Open \\
    Knowledge-limit failure & Model lacks the parametric knowledge needed to answer without a supplied passage. & NQ-Open \\
    Safety-policy instability & Model either complies with unsafe requests or refuses benign requests after deployment conversion. & Safety suite \\
    \bottomrule
  \end{tabular}
\end{table*}

This taxonomy guided the interpretation of the main results. SQuAD P3 is mainly
a task-contract failure; NQ-Open at E2--E4 combines knowledge limits with
runtime-interface failure; HaluEval's initial low scores were primarily a parser
and runtime-interface issue; and safety variation reflects policy instability
rather than a monotonic quantization effect.

\subsection{Why MT-Bench Is Not in the 4-Tier Heatmap}

MT-Bench is a classic eight-category, open-ended, judge-scored benchmark. Its
unit of evaluation is an 80-question conversation set rather than a 200-example
reference-scored sample. We therefore do not mix it into the P0/P1/P2/P3
reference-task heatmaps. Instead, we use MT-Bench to test a different question:
whether informal prompting changes open-ended instruction-following quality and
whether model-family rankings are stable under deployment conversion. The main
paper reports setting-level and family-level P0/P3 deltas; category-level radar
plots are included in the released analysis notebook.

\subsection{Interpreting Missing Cells}

Missing cells in EdgeBench are not silently imputed. They arise from three
sources: (1) model families that require a different runtime, such as Falcon-E
with \texttt{bitnet.cpp}; (2) tokenizer or chat-template incompatibilities that
prevent reliable generation for a task/setting; and (3) optional baseline models
with partial coverage. We report aggregates over available cells and avoid
claiming a comparison when the matched model--setting--prompt set is incomplete.
For prompt-sensitivity claims, we use matched P0/P3 pairs wherever possible.

\subsection{Reproducibility Notes}

The released repository contains three layers of artefacts: raw generations,
per-run score files, and aggregate CSV summaries. The intended workflow is:
(1) run generation jobs with a model key, setting key, dataset key, and prompt
tier; (2) score raw outputs with the task-specific evaluator; (3) aggregate
per-run metrics into \texttt{results/summary.csv}; and (4) derive paper tables
and figures from the summary plus MT-Bench \texttt{scores.jsonl} files. This
separation is deliberate. It allows parser fixes, such as the SQuAD EOS
correction and HaluEval first-line correction, to be applied without rerunning
model inference.

For a new model, the minimal experiment matrix is P0/P3 for SQuAD v2,
TruthfulQA-MC1, NQ-Open, and MT-Bench under E0 plus the intended edge setting.
The full EdgeBench matrix then expands to all four prompt tiers for
reference-scored tasks and all five settings for deployment analysis. We
recommend running smoke tests with a small example limit before launching the
full 200-example runs, because runtime failures are common in the edge setting
and should be discovered before queueing large jobs.

\subsection{Threats to Validity}

\paragraph{Prompt realism.}
P1--P3 are hand-authored diagnostic prompts. They represent plausible levels of
prompt specificity, but they are not drawn from a user study. The claims should
therefore be read as robustness-to-controlled-prompt-degradation claims, not as
population-level claims about all novice users.

\paragraph{Hardware realism.}
E2--E4 are CPU-only GGUF deployments with constrained thread and memory settings.
They approximate a local edge environment, but they do not cover phones,
NPUs, mobile GPUs, or vendor-specific acceleration. We removed phone-device
comparisons from the current paper because no phone measurements were collected.

\paragraph{Judge reliability.}
MT-Bench relies on GPT-4o-mini as an automatic judge. The judge is useful for
relative instruction-following comparisons, but it introduces model-specific
judgment noise. For this reason, MT-Bench conclusions are reported separately
from reference-scored QA and classification tasks.

\paragraph{Safety scope.}
The safety suite measures refusal/compliance patterns with a keyword classifier.
It is appropriate for detecting deployment instability, but insufficient for
policy compliance claims. A future version should include human annotation or a
stronger policy judge.

\end{document}
